\documentclass{beamer}
\usepackage[utf8]{inputenc}

\usetheme{Madrid}
\usecolortheme{default}

% Informações da capa
\title{South Korea}
\subtitle{Language, Culture, and the Job Market}
\author{Pedro Henrique Peixoto Cardoso}
\date{October 2026}

\begin{document}

% Slide 1: Capa
\begin{frame}
    \titlepage
\end{frame}

\begin{frame}
    \frametitle{How the English Language Arrived}
    \begin{itemize}
        \item \textbf{Missionaries:} Protestant missionaries introduced English alongside Western curricula(1885);
        \item \textbf{Post-WWII US Presence:} Expansion driven by US military and financial aid after WWII and the Korean War (1950–1953);
        \item \textbf{Social Status:} In the 70s and 80s, English mastery became a status symbol and a requirement for professional success;
        \item \textbf{The Konglish Phenomenon:} A linguistic fusion where English words are adapted to the Korean alphabet (Hangul).
    \end{itemize}
\end{frame}

\begin{frame}
    \frametitle{Other Languages Spoken}
    \begin{itemize}
        \item \textbf{Korean (Hangul):} The official language and the main symbol of national pride;
        \item \textbf{English:} The most widely studied foreign language, used in media, trade, and tech;
        \item \textbf{Historical Context:} Classical Chinese was used by the nobility before 1446. Japanese terms were used during the occupation (1910–1945) but later removed to preserve the Korean language.
    \end{itemize}
\end{frame}

\begin{frame}
    \frametitle{Job Market}
    \begin{itemize}
        \item \textbf{Corporate Culture:} Highly competitive, long hours, and dominated by \textit{Chaebols} (family-owned conglomerates);
        \item \textbf{Social Challenges:} Women advanced in the workforce, but clashes with traditional patriarchy created the \textit{Sampo Generation} (youth giving up dating, marriage, and children);
        \item \textbf{Demographic Crisis:} A low fertility rate (0.72) and an aging population pose major risks to the future workforce.
    \end{itemize}
\end{frame}

\begin{frame}
    \frametitle{Brazilians in South Korea}
    \begin{itemize}
        \item \textbf{Professional Expatriation:} Brazilian executives and professionals sent by multinational companies face challenges adapting to the corporate culture;
        \item \textbf{Tourism and Study:} A growing number of Brazilians travel to South Korea to learn Hangul and explore the country, heavily motivated by the global appeal of K-pop and K-dramas.
    \end{itemize}
\end{frame}

\begin{frame}
    \frametitle{General Curiosities}
    \begin{itemize}
        \item \textbf{Hangul Day (Oct 9):} A national holiday celebrating the alphabet created by King Sejong in 1443;
        \item \textbf{The Suneung Exam:} A brutal 8-hour university entrance exam that deeply impacts a student's future career and social standing;
        \item \textbf{Pop Tourism:} Landmarks like the Han River, Namsan Tower, and K-drama filming sets are massive global attractions;
        \item \textbf{Konglish Vocab:} English words adapted to Korean phonology (e.g., \textit{handphone} for cell phone, \textit{ice cream}).
    \end{itemize}

\end{frame}
\begin{frame}
    \centering
    \Huge \textbf{Thank you for your attention!}
\end{frame}


\begin{frame}[allowframebreaks]
    \frametitle{References}
    
    \footnotesize 
    
    \begin{thebibliography}{99}
    
    \bibitem{araujo2024}
    ARAÚJO, I. I. S.; MOREIRA, J. B. D.
    \newblock \textit{Korean Wave: A influência do marketing de entretenimento na propagação da cultura sul-coreana}.
    \newblock Revista Multidisciplinar do Centro Universitário FIPMoc, Montes Claros, v. 37, n. 1, p. 86-121, 2024.
    
    \bibitem{coelho2019}
    COELHO, A. R.
    \newblock \textit{Expatriados: Expectativas, percepções e realidade. Um estudo sobre as experiências de expatriados brasileiros em países asiáticos}.
    \newblock Artigo (Mestrado Profissional em Administração) – FECAP, São Paulo, 2019.
    
    \bibitem{marino2020}
    MARINO, C. B. G.
    \newblock \textit{O fenômeno Konglish na música popular coreana (K-Pop)}.
    \newblock Monografia (Bacharelado em Tradução) – UFU, Uberlândia, 2020.
    
    \bibitem{monteiro2025}
    MONTEIRO, M. E. M.
    \newblock \textit{Crise demográfica na Coreia do Sul: Causas e consequências da queda na taxa de natalidade para a economia sul-coreana}.
    \newblock TCC (Relações Internacionais) – PUC-Campinas, Campinas, 2025.
    
    \end{thebibliography}
\end{frame}

\end{document}
